% Generated draft: requires mathematical and source review.
\documentclass[11pt]{article}
\usepackage[margin=25mm]{geometry}
\usepackage{fontspec}
\setmainfont{Latin Modern Roman}
\usepackage{amsmath,amssymb,mathtools,unicode-math}
\setmathfont{Latin Modern Math}
\usepackage{enumitem,needspace,xurl,hyperref}
\setlength{\parindent}{0pt}
\setlength{\parskip}{5pt}
\setlength{\emergencystretch}{3em}
\newcommand{\N}{\mathbb{N}}
\newcommand{\Z}{\mathbb{Z}}
\newcommand{\Q}{\mathbb{Q}}
\newcommand{\R}{\mathbb{R}}
\newcommand{\C}{\mathbb{C}}
\newcommand{\F}{\mathbb{F}}
\newcommand{\A}{\mathbb{A}}
\newcommand{\PP}{\mathbb{P}}
\newcommand{\E}{\mathbb{E}}
\newcommand{\eps}{\varepsilon}
\newcommand{\dd}{\,\mathrm{d}}
\begin{document}
\section*{UnsolvedMath ID: 30004936}
\textbf{Source problem code:} catalogue entry
\subsection{Definitions and mathematical statement}
The supplied record for permanent UnsolvedMath ID \(30004936\) contains a workshop abstract, not a mathematical problem statement. Its recorded statement status is \texttt{unrecoverable}, and no clean statement is supplied. Consequently, no precise conjecture, theorem, or quantified mathematical question can be recovered from this evidence.

The abstract describes connections between cluster algebras, scattering amplitudes, and positive geometries, including amplituhedra. It attributes the development of cluster algebras to Fomin and Zelevinsky, approximately twenty years before the abstract, with initial motivations in total positivity and canonical bases. It also describes a program begun by Arkani-Hamed and collaborators approximately ten years before the abstract to reformulate the foundations of quantum field theory, motivated by seeking deeper mathematical principles underlying quantum mechanics and spacetime. According to the abstract, this approach encodes particle scattering amplitudes in geometric objects and involves connections with cluster algebras. These are supplied contextual descriptions, not a specified mathematical assertion to be proved.

The stated workshop objective is to explore connections among these subjects. The passage supplies no definitions of the relevant cluster algebras, scattering amplitudes, positive geometries, or amplituhedra; it specifies no particular objects, hypotheses, domains, boundary conditions, quantifiers, or target correspondence. Introducing any such choices would create a new problem rather than recover the supplied one. A mathematical statement therefore remains unavailable pending identification of an actual problem passage and its necessary definitions and hypotheses.

The registry labels the record \texttt{open}. Its supplied verification note identifies the passage as a workshop abstract and says that no separate clean statement is promoted, while the original record remains publicly visible for provenance. The supplied dated literature assessment likewise treats the extracted claim as unresolved because it is not a definite conjecture or theorem. Neither label establishes the open status of a specified mathematical proposition. The abstract lists Mathematics Subject Classification (2020) codes 81U10, 13F60, 14M15, and 51M20.
\subsection{Short English statement}
This record describes a workshop on connections between cluster algebras, scattering amplitudes, and positive geometries. It contains no definite mathematical problem to formalize. Its supplied \texttt{open} label and \texttt{unrecoverable} statement status are retained as registry qualifications; a precise statement requires an identifiable problem passage with definitions and hypotheses.
\subsection{Sources}
\begin{itemize}
\item Local UnsolvedMath registry: \texttt{data/unsolvedmath/problems.json}, with \texttt{catalog/problems.json} as the additional-record catalog; supplied permanent ID \(30004936\), exact source problem code \texttt{catalogue entry}, title \emph{Cluster Algebras Scattering Amplitudes and Positive Geometries}, in the Oberwolfach Reports Open Problems collection. This reference identifies the supplied registry record; the local files were not independently inspected.
\item Supplied source citation: \emph{Mini-Workshop: Scattering Amplitudes, Cluster Algebras, and Positive Geometries} (2022), Oberwolfach Reports, DOI 10.4171/owr/2021/57, pp.\ 3101--3133. Supplied source URL and literature-source link labeled ``original Oberwolfach report'': \url{https://doi.org/10.4171/owr/2021/57}. The record separately gives proposed year 2021 and the tag \texttt{owr-year:2021}; their relationship to the citation year 2022 is not established by the supplied evidence.
\item Attribution in the supplied abstract: Fomin and Zelevinsky are credited with developing cluster algebras; Arkani-Hamed and collaborators are associated with the described quantum-field-theory program. No individual problem proposer is supplied.
\item Supplied literature assessment dated exactly 2026-08-22: status \texttt{open}; the extracted passage is assessed as workshop motivation rather than a definite mathematical conjecture or theorem, so the claim is treated as unresolved and not upgraded beyond open. This is supplied dated curation, not independently verified. The background describes it as literature triage derived from supplied OWR curation and calls for independent verification of time-sensitive status and cited claims. The source link above has not been visited for this draft.
\end{itemize}
\end{document}
